\section{Cross-sample prediction churn}
\label{sec:preliminaries}
\label{sec:measurement}

\paragraph{Setting.}
A learner $\mathcal{A}$ maps a training set
$S = \{(x_i, y_i)\}_{i=1}^{N}$ drawn iid from population $\mathcal{D}$
to a classifier $f_S = \mathcal{A}(S)$ that returns a class
$\hat{y}(x) = \arg\max_c f_S(x)_c$ at each test example $x$.
Cross-sample churn quantifies how $\hat{y}(x)$ changes when $S$
is replaced by an independent iid draw from $\mathcal{D}$.

\paragraph{Definition.}
Given two iid training samples $S_A, S_B \sim \mathcal{D}^N$ and a
test example $x$, the \emph{cross-sample churn} of $\mathcal{A}$ at
$x$ is
\begin{equation}
  \rho(\mathcal{A}, x) =
  \mathbb{P}_{S_A, S_B}\!\big[\arg\max f_{S_A}(x) \neq \arg\max f_{S_B}(x)\big].
\end{equation}
A churn-aware report supplements the standard accuracy with the
expected per-example argmax disagreement
$\bar\rho(\mathcal{A}) = \mathbb{E}_x [\rho(\mathcal{A}, x)]$, and
optionally with a distributional analogue
$\bar\rho_{\mathrm{KL}}(\mathcal{A})
= \mathbb{E}_x \big[\tfrac{1}{2}(\mathrm{KL}(f_{S_A}(x) \| f_{S_B}(x))
                                + \mathrm{KL}(f_{S_B}(x) \| f_{S_A}(x)))\big]$.
We use ``cross-sample churn'' (or just ``churn'' when context is
clear) for $\bar\rho$ and ``class-flip rate'' for the same quantity
reported as a percentage; the two names refer to the same quantity.

\paragraph{Per-example, not aggregate.}
Every aggregate metric (accuracy, AUROC, precision, recall, $F_1$,
AP) is a population average that summarises out per-example variation;
churn counts disagreements before any averaging.  We pre-screen
to ERM accuracy $\geq$ majority-class baseline $+\preregFilterPP$pp
so the gap we measure is learned-decision performance rather than
majority shuffling.  On every dataset, the per-example class-flip
rate exceeds every aggregate-metric drift, typically by
$2\text{--}10\times$ (e.g.\ cyp2d6\_substrate:
$|\Delta\text{recall}|=\cypTwoDeltaRecall$\,pp,
$|\Delta\text{acc}|=\cypTwoDeltaAcc$\,pp;
Appendix~\ref{app:additional_metrics}).  The disagreement also
concentrates on the minority class on imbalanced datasets
($2\text{--}4\times$ majority on BBB-Martins, BBBP, CYP2D6-Sub;
Appendix~\ref{app:per_class_churn}) --- the class that drives
next-step decisions.

\paragraph{Churn correlates with task difficulty, but is not
explained by it.}
At the dataset level, the mean class-flip rate correlates with mean
ERM accuracy at Pearson $\churnAccCorrPearson$ (Spearman
$\churnAccCorrSpearman$) across the \nDatasetsHeadline{} chemistry
benchmarks --- harder tasks show more churn.  At the
per-method level, decoupling: deep ensembles raise mean accuracy by
$\deepEnsAccDelta$\,pp without reducing churn, while twin-bootstrap
cuts churn a median $\twinMedianReduction\%$ at $\twinAccDelta$\,pp
mean accuracy cost (Table~\ref{tab:main}).  Churn is not a
recoding of accuracy.

\paragraph{Bootstrap estimator.}
A practitioner has one observed training set $S$, not a population
$\mathcal{D}$.  We estimate $\rho$ by drawing two independent
bootstraps from $S$ (size $N$, with replacement); each bootstrap is
a draw from the empirical distribution $\hat{\mathcal{D}}_S$, which
converges to $\mathcal{D}$ as $N \to \infty$.  All
numerical magnitudes we report are with respect to this bootstrap
proxy --- the quantity computable from a single observed training
set.

\paragraph{Measurement protocol.}
A single fixed seed --- which we call the \emph{canonical} seed
--- determines the train/test split once and for all; the resulting
training pool and test set are reused across every method and every
train seed, so test-side variance is zero by construction within a
canonical seed.  For each train seed $s \in \{1, \dots, \nSeeds\}$ we
draw a bootstrap $S^{(s)}$ from the canonical training pool and train
$f^{(s)} = \mathcal{A}(S^{(s)})$ on it.  Cross-sample churn is the average of
$\mathbf{1}[\arg\max f^{(s)}(x) \neq \arg\max f^{(s')}(x)]$ over all
$\binom{\nSeeds}{2}{=}\nSeedPairs$ pairs $(s, s')$ and over the canonical test set.
Confidence intervals come from $10{,}000$-sample bootstrap on the
seed-pair distribution; cross-method deltas are paired.

We replicate the entire protocol on $\nCanonicalSeeds$ independent
canonical seeds ($\canonicalSeed$, $7$, $42$); magnitudes quoted in
prose are across-seed means and the $\lambda$-selection rule picks
$\lambda{=}300$ on every seed independently
(Appendix~\ref{app:seed_sensitivity}).  Pretrained-backbone
cross-checks (GIN, ChemBERTa, Waterbirds) are single-seed.

\paragraph{Datasets.}
\label{sec:datasets}
Seventeen chemistry benchmarks ($N \in [\,304, 4658\,]$):
MoleculeNet~\citep{Wu2018} (BACE, BBBP, ClinTox),
the TDC~\citep{huang2021therapeutics} ADME and Tox suites
(HIA\_Hou, Bioavailability\_Ma, Pgp\_Broccatelli, BBB\_Martins,
CYP\{2C9,2D6,3A4\}-Sub, hERG, DILI, AMES, Skin\_Reaction), and
materials-science benchmarks (TADF~\citep{jablonka2026clever, Huang2024},
MOF-thermal-stability, MOF-solvent-removal~\citep{jablonka2026clever, Nandy2022}).
A pre-registered health filter
(ERM accuracy $\geq$ majority $+ \preregFilterPP$pp on a canonical
test set $\geq \preregMinTestSize$ examples) keeps nine headline
datasets, flags three borderline datasets, and excludes five; per-dataset
filter outcomes are in Appendix~\ref{app:repro}.  Featurisation is
the standard chemistry choice per dataset family ($2048$-bit Morgan
radius-2 fingerprints for MoleculeNet/TDC; $+217$ RDKit descriptors
for TADF; $174$ RAC descriptors plus GSA for
MOFs~\citep{Moosavi2020}; full details in
Appendix~\ref{app:repro}).  The default architecture is a
$256$-unit MLP trained from scratch on these features.

